\PageHeading{01}{项目概览}{面向门店销售员的、由产品证据驱动的营销内容工作台}
\begin{note}
\Lead{项目定位}\quad 知序将“理解客户、查找资料、共创大纲、生成图文、核验与导出”组织成一个可追溯工作流。ima 提供产品事实与原图，自研 Agent 负责业务判断、内容策略、结构控制与审核。公开大模型是其中的写作工具。
\end{note}
\Stat{8}{类场景}{来自赛题建议场景目录}\hfill
\Stat{4}{个渠道}{公众号、小红书、朋友圈、海报}\hfill
\Stat{89}{项测试}{已交付基线：0失败、0跳过}
\Sub{要解决的业务问题}
销售员需要把同一产品面向不同客户讲清楚。难点同时存在于三个环节：资料版本和型号容易混淆；产品参数需要转化为具体场景的选购问题；不同平台对标题、结构、配图与行动引导的要求不同。知序将这些规则放入流程和数据结构，减少对单条自由提示词的依赖。
\Sub{本项目的核心设计}
\begin{enumerate}
\item \Lead{先确认方向，再展开正文。}销售员可手动编辑或让模型修改大纲；保存产生新版本，确认后才能启动新流程正文生成。
\item \Lead{事实、图片与型号一起流转。}事实附带原文片段、限制条件和来源记录；图片绑定产品及位置，渲染只读取已有任务中的受控素材。
\item \Lead{把审核保留到交付。}规则检查与独立审稿调用共同形成审核结果；必要时仅返修一次。待确认项随成品保留，阻断稿只提供带标记预览。
\item \Lead{从文本方案走向实际文件。}四渠道已生成 HTML、PNG 和来源 ZIP；正文无法完整放入图版时明确报错，避免截断条件。
\end{enumerate}
\begin{attention}
\Lead{当前完成状态}\quad 本机七步流程、四渠道排版及真实模型调用已有运行证据。本文三组案例使用获授权的\textbf{本地资料}；ima 连接器已实现，但当前额度限制下尚未完成新流程的 ima 端到端验收。本地样例不能替代赛事对 ima 事实和图片来源的正式核验。
\end{attention}
\Sub{阅读导航}
\begin{tabularx}{\linewidth}{P{34mm}Y}
\toprule
\Lead{评委快速了解} & 赛题对照（第\pageref{sec:02}页）、功能流程（第\pageref{sec:03}页）、三组结果（第\pageref{sec:08}--\pageref{sec:10}页）。\\
\Lead{技术方案评审} & 架构、证据链与 Agent（第\pageref{sec:04}--\pageref{sec:06}页），验证与部署（第\pageref{sec:11}--\pageref{sec:12}页）。\\
\Lead{提交前核对} & 限制与迭代（第\pageref{sec:13}页）、材料清单（第\pageref{sec:14}页）、提示词原文（第\pageref{sec:15}页）。\\
\bottomrule
\end{tabularx}
\SourceNote{手册第3--5页；已交付实现8574709及接口基线163616。制作期间的工作区增量单独标注，未并入本报告完整流程验收。}
